Dels esdeveniments $A$ i $B$ d'un experiment aleatori se sap que $P(A)=0{,}5$, $P(B)=0{,}4$ i
$P\left(\overline{A}\cup\overline{B}\right)=0{,}8$.

\begin{apartats}

\begin{tria}{operacions-tasca}
\itemtria{original}{0,75}{1,25}
Calcula $P(A\cap B)$ i $P(A\cup B)$.

\begin{solucio}
Per la llei de De Morgan, $\overline{A}\cup\overline{B}=\overline{A\cap B}$, i per tant
$P(A\cap B)=1-0{,}8=0{,}2$.\\
$P(A\cup B)=0{,}5+0{,}4-0{,}2=0{,}7$.
\end{solucio}

\itemtria{fites}{0,75}{1,25}
Poden ser incompatibles dos esdeveniments $C$ i $D$ amb $P(C)=0{,}7$ i $P(D)=0{,}6$? Quin és el valor més
petit que pot tenir $P(C\cap D)$?

\begin{solucio}
Si fossin incompatibles, seria $P(C\cup D)=0{,}7+0{,}6=1{,}3>1$, cosa impossible: \textbf{no} ho poden
ser.\\
Com que $P(C\cup D)\le1$, $P(C\cap D)=P(C)+P(D)-P(C\cup D)\ge0{,}7+0{,}6-1=0{,}3$. El valor més petit és
$0{,}3$, quan $C\cup D$ és l'esdeveniment segur.
\end{solucio}
\end{tria}

\apartat[1,25]{1}
Calcula $P\left(A\cap\overline{B}\right)$ i $P\left(\overline{A}\cap\overline{B}\right)$.

\begin{solucio}
$P\left(A\cap\overline{B}\right)=P(A)-P(A\cap B)=0{,}5-0{,}2=0{,}3$.\\
$P\left(\overline{A}\cap\overline{B}\right)=P\left(\overline{A\cup B}\right)=1-0{,}7=0{,}3$.
\end{solucio}

\begin{nomesllarg}
\apartat{0,75}
Són incompatibles $A$ i $B$? I independents? Justifica-ho.

\begin{solucio}
\textbf{No} són incompatibles: $P(A\cap B)=0{,}2\neq0$.\\
\textbf{Sí} que són independents: $P(A)\cdot P(B)=0{,}5\cdot0{,}4=0{,}2=P(A\cap B)$.
\end{solucio}
\end{nomesllarg}

\end{apartats}
